%% toc: P3 28 28.1,28.2,28.3,28.4
\chapter{Quantifier Elimination, Existential Closure, and Model Companions}
\label{sec:p03-w03-p03-020-qe}

Quantifier elimination replaces every formula, modulo a theory, by a formula
without quantifiers. It is the strongest of four related properties that this
chapter keeps apart: quantifier elimination for a theory, model completeness
of a theory, existential closure of a model in a class of models, and the
relations of model companion and model completion between two theories. The
inputs throughout are a first-order signature \(\Lang\) and an
\(\Lang\)-theory \(T\); the conclusions are equivalences of formulas modulo
\(T\), elementarity of embeddings between models of \(T\), and completeness of
\(T\) together with the atomic diagram of a structure.

The chapter supplies the proofs that the earlier entries state or outline.
The definition of quantifier elimination and the one-quantifier method are in
\cref{sec:qe-categoricity-01}; the definitions of existentially closed models,
model companions and model completions are
\cref{def:p03-existentially-closed,def:p03-model-companion}, and the cited
statements of Robinson's test and of the field examples are
\cref{thm:p03-robinson-model-completeness,thm:p03-field-model-completions}.
Below, the semantic and syntactic criteria for elimination are proved, Robinson's
test is proved together with its diagram and existential-normal-form forms, the
basic existence, uniqueness and embedding facts for companions are proved,
dense linear orders and their endpoint variants are treated by back-and-forth,
and algebraically closed fields of fixed characteristic are shown to be the
model completion of fields of that characteristic, including the algebra of
algebraic closures. In the ordered-ring language, real closures of ordered
fields are constructed and real closed fields are shown to be the model
completion of ordered fields, with completeness over the diagram of each
ordered field. The effective elimination procedures and their use for
decision problems are in \cref{sec:p03-w03-p03-022-decidability}. The general
criteria follow the treatment in Kruckman's lecture notes
\cite[sec.~2.2--2.3, pp.~16--22]{Kruckman2020}; the algebraic lemmas, the
companion facts and the back-and-forth arguments are proved here in full.

\section{Semantic and syntactic criteria for quantifier elimination}
\label{sec:p03-w03-p03-020-criteria}

\begin{hbconvention}{def:p03-w03-p03-020-standing}{Standing conventions}
Structures are nonempty. Quantifier-free formulas include \(\top\) and
\(\bot\). A \emph{literal} is an atomic formula or the negation of one, and a
\emph{primitive formula} is one of the form \(\exists y\,\theta\) with
\(\theta\) a finite conjunction of literals. For a formula
\(\varphi(\bar x)\), the tuple \(\bar x=(x_1,\ldots,x_n)\) lists distinct
variables including all free variables of \(\varphi\); \(n=0\) is allowed.
Two formulas \(\varphi(\bar x)\) and \(\psi(\bar x)\) are
\emph{\(T\)-equivalent} when
\(T\entails\forall\bar x\,(\varphi\leftrightarrow\psi)\).
For an \(\Lang\)-structure \(\mathcal M\) and \(\bar a\in M^n\), the
\emph{quantifier-free type}
\(\operatorname{qftp}^{\mathcal M}(\bar a)\) is the set of quantifier-free
formulas \(\psi(\bar x)\) with \(\mathcal M\entails\psi(\bar a)\). For
\(n=0\) it is the set of quantifier-free sentences true in \(\mathcal M\).
An \emph{existential} formula is \(\exists\bar y\,\theta\) and a
\emph{universal} formula is \(\forall\bar y\,\theta\), with \(\theta\)
quantifier-free and \(\bar y\) possibly empty. \(T_\forall\) is the set of
universal sentences entailed by \(T\), and \(T_{\forall\exists}\) the set of
entailed sentences of the form \(\forall\bar x\,\exists\bar y\,\theta\) with
\(\theta\) quantifier-free. Diagrams and the parameter signature
\(\Lang(A)\) are as in \cref{def:parameter-expansion,def:diagrams}.
\end{hbconvention}

Quantifier-free types replace the common-substructure formulation of the
criteria. The two agree when the tuple is nonempty or the signature has
constant symbols, because \(\bar a\) and \(\bar b\) have the same
quantifier-free type exactly when \(t(\bar a)\mapsto t(\bar b)\) is a
well-defined isomorphism between the substructures they generate. The type
formulation also covers sentences in a signature without constants, where the
generated substructure would be empty.

\begin{hbtheorem}{thm:p03-w03-p03-020-primitive}{proved=annex:p03-w03-p03-020-primitive}{Syntactic criterion: primitive formulas suffice}
A theory \(T\) eliminates quantifiers if and only if every primitive formula
\(\exists y\,\theta(\bar x,y)\) is \(T\)-equivalent to a quantifier-free
formula \(\psi(\bar x)\).
\end{hbtheorem}

This is the justification of the one-quantifier method of
\cref{sec:qe-categoricity-01}. The algorithmic version, which also keeps
track of computability, is \cref{thm:p03-w03-p03-022-one-step}.

\begin{hbtheorem}{thm:p03-w03-p03-020-semantic}{proved=annex:p03-w03-p03-020-semantic}{Semantic criterion for one formula}
Let \(T\) be a theory and \(\varphi(\bar x)\) a formula. The following are
equivalent.
\begin{enumerate}
\item \(\varphi(\bar x)\) is \(T\)-equivalent to a quantifier-free formula
  \(\psi(\bar x)\).
\item For all models \(\mathcal M,\mathcal N\) of \(T\) and all
  \(\bar a\in M^n\), \(\bar b\in N^n\) with
  \(\operatorname{qftp}^{\mathcal M}(\bar a)=\operatorname{qftp}^{\mathcal N}(\bar b)\),
  if \(\mathcal M\entails\varphi(\bar a)\) then
  \(\mathcal N\entails\varphi(\bar b)\).
\end{enumerate}
\end{hbtheorem}

\begin{hbtheorem}{thm:p03-w03-p03-020-one-point}{proved=annex:p03-w03-p03-020-one-point}{One-point extension test}
A theory \(T\) eliminates quantifiers if and only if the following holds:
whenever \(\mathcal M,\mathcal N\entails T\),
\(\bar a\in M^n\) and \(\bar b\in N^n\) have the same quantifier-free type,
and \(c\in M\), there are an elementary extension
\(\mathcal N\preccurlyeq\mathcal N'\) and \(d\in N'\) such that
\(\operatorname{qftp}^{\mathcal M}(\bar a,c)=\operatorname{qftp}^{\mathcal N'}(\bar b,d)\).
\end{hbtheorem}

The test is the form used for fields below. It asks only that a single new
element be matched, and the elementary extension may be chosen to contain the
match; the price is that the method gives no explicit equivalent formula.

\begin{hbcorollary}{thm:p03-w03-p03-020-qe-consequences}{proved=annex:p03-w03-p03-020-qe-consequences}{What elimination gives}
Let \(T\) eliminate quantifiers in \(\Lang\).
\begin{enumerate}
\item If \(T\subseteq T'\) for an \(\Lang\)-theory \(T'\), then \(T'\)
  eliminates quantifiers.
\item If \(\Lang'\) is \(\Lang\) with new constant symbols, then \(T\),
  read as an \(\Lang'\)-theory, eliminates quantifiers in \(\Lang'\).
\item \(T\) is model complete.
\item For every \(\Lang\)-structure \(\mathcal A\), the
  \(\Lang(A)\)-theory \(T\cup\operatorname{Diag}(\mathcal A)\) is complete:
  for each \(\Lang(A)\)-sentence \(\sigma\), it entails \(\sigma\) or
  \(\neg\sigma\). It is consistent exactly when \(\mathcal A\) embeds in a
  model of \(T\).
\item If \(T\) is consistent and any two models of \(T\) satisfy the same
  atomic sentences, then \(T\) is complete.
\end{enumerate}
\end{hbcorollary}

Item~4 is the diagram form of elimination. It is the property that turns a
companion with elimination into a model completion in
\cref{thm:p03-w03-p03-020-qe-completion}.

\subsection{Model completeness and Robinson's test}
\label{sec:p03-w03-p03-020-robinson}

Model completeness is defined in \cref{def:theory-completeness}. For
\(\mathcal M\subseteq\mathcal N\), say that \(\mathcal M\) is
\emph{existentially closed in} \(\mathcal N\) when every existential
\(\Lang(M)\)-sentence true in \(\mathcal N_M\) is true in \(\mathcal M_M\).
\Cref{def:p03-existentially-closed} quantifies this over all extensions that
are models of a given theory.

\begin{hblemma}{thm:p03-w03-p03-020-existential-preservation}{proved=annex:p03-w03-p03-020-existential-preservation}{Existential normal form from preservation}
Let \(T\) be a theory and \(\varphi(\bar x)\) a formula. Then
\(\varphi\) is \(T\)-equivalent to an existential formula if and only if,
for all models \(\mathcal M\subseteq\mathcal N\) of \(T\) and all
\(\bar a\in M^n\), \(\mathcal M\entails\varphi(\bar a)\) implies
\(\mathcal N\entails\varphi(\bar a)\). Dually, \(\varphi\) is
\(T\)-equivalent to a universal formula if and only if truth of
\(\varphi\) passes from \(\mathcal N\) down to every submodel
\(\mathcal M\) of \(T\).
\end{hblemma}

\begin{hbtheorem}{thm:p03-w03-p03-020-robinson}{proved=annex:p03-w03-p03-020-robinson}{Robinson's test, with its diagram and normal-form bridges}
For a theory \(T\) the following are equivalent.
\begin{enumerate}
\item \(T\) is model complete.
\item Whenever \(\mathcal M\subseteq\mathcal N\) are models of \(T\),
  \(\mathcal M\) is existentially closed in \(\mathcal N\).
\item Every formula is \(T\)-equivalent to an existential formula.
\item Every formula is \(T\)-equivalent to a universal formula.
\item For every model \(\mathcal M\) of \(T\), the \(\Lang(M)\)-theory
  \(T\cup\operatorname{Diag}(\mathcal M)\) is complete; equivalently, it
  entails \(\operatorname{ElDiag}(\mathcal M)\).
\end{enumerate}
\end{hbtheorem}

The step from 2 to 3 is the existential-normal-form bridge: it passes from a
property of pairs of models to a syntactic normal form through
\cref{thm:p03-w03-p03-020-existential-preservation}. The step from 3 to 1
is the reverse bridge. Items 1 and 5 are the diagram bridge in both
directions. This proves the statement recorded with a citation at
\cref{thm:p03-robinson-model-completeness}, and adds items 4 and 5.

\begin{hblemma}{thm:p03-w03-p03-020-ec-elementary}{proved=annex:p03-w03-p03-020-ec-elementary}{Existential closure and elementary extensions}
If \(\mathcal A\subseteq\mathcal B\) and \(\mathcal A\) is existentially
closed in \(\mathcal B\), then there is an elementary extension
\(\mathcal A\preccurlyeq\mathcal A'\) and an embedding of \(\mathcal B\)
into \(\mathcal A'\) that is the identity on \(A\).
\end{hblemma}

\begin{hbtheorem}{thm:p03-w03-p03-020-ae}{proved=annex:p03-w03-p03-020-ae}{Model complete theories are \(\forall\exists\)-axiomatized}
If \(T\) is model complete, then \(T_{\forall\exists}\) and \(T\) have the
same models.
\end{hbtheorem}

\subsection{Companions and completions}
\label{sec:p03-w03-p03-020-companions}

\begin{hblemma}{thm:p03-w03-p03-020-companion-embedding}{proved=annex:p03-w03-p03-020-companion-embedding}{Companions by mutual embedding}
Let \(T\) and \(T^*\) be theories in the same signature. Then
\(T_\forall=T^*_\forall\) if and only if every model of \(T\) embeds in a
model of \(T^*\) and every model of \(T^*\) embeds in a model of \(T\).
\end{hblemma}

\begin{hbtheorem}{thm:p03-w03-p03-020-companion-unique}{proved=annex:p03-w03-p03-020-companion-unique}{Uniqueness of model companions}
If \(T_1^*\) and \(T_2^*\) are model companions of \(T\), then they have
the same models. In particular a model completion of \(T\), when it exists,
is unique up to having the same models.
\end{hbtheorem}

\begin{hbtheorem}{thm:p03-w03-p03-020-companion-ec}{proved=annex:p03-w03-p03-020-companion-ec}{Models of a model companion are the existentially closed models}
Let \(T^*\) be a model companion of \(T\).
\begin{enumerate}
\item Every model of \(T^*\) is existentially closed in each of its
  extensions that is a model of \(T_\forall\), in particular in each
  extension that is a model of \(T\).
\item If \(T\) has the same models as \(T_{\forall\exists}\), then the
  models of \(T^*\) are exactly the models of \(T\) that are existentially
  closed among models of \(T\) in the sense of
  \cref{def:p03-existentially-closed}.
\end{enumerate}
\end{hbtheorem}

\begin{hbtheorem}{thm:p03-w03-p03-020-qe-completion}{proved=annex:p03-w03-p03-020-qe-completion}{A companion with elimination is a model completion}
Let \(T^*\) satisfy \(T^*_\forall=T_\forall\) and eliminate quantifiers.
Then \(T^*\) is the model completion of \(T\). Explicitly, for every model
\(\mathcal M\) of \(T\), the \(\Lang(M)\)-theory
\(T^*\cup\operatorname{Diag}(\mathcal M)\) is consistent and complete.
\end{hbtheorem}

\begin{hbwarning}{def:p03-w03-p03-020-distinctions}{Four notions that are not quantifier elimination}
\begin{enumerate}
\item \emph{Existential closure} is a property of one model relative to a
  class of extensions. It says nothing about formulas with universal
  quantifiers, and it does not make any theory eliminate quantifiers. A model
  of \(\mathsf{DLO}\) is existentially closed among linear orders by
  \cref{thm:p03-w03-p03-020-companion-ec,thm:p03-w03-p03-020-dense-variants},
  yet the theory of linear orders has no elimination: the formula
  \(\exists y\,(x<y)\) is not equivalent to a quantifier-free formula in
  \(x\) in the two-element order, as noted after
  \cref{thm:p03-w03-p03-022-dlo-step}.
\item \emph{Model completeness} is weaker than elimination. By
  \cref{thm:p03-w03-p03-020-qe-consequences} elimination implies it, and
  \cref{thm:p03-w03-p03-020-robinson} shows that it amounts to elimination
  down to existential, not quantifier-free, formulas.
  \Cref{thm:p03-w03-p03-020-marked-point} gives a model complete theory
  without elimination.
\item A \emph{model companion} is a theory related to another theory. Its
  existence and uniqueness (\cref{thm:p03-w03-p03-020-companion-unique}) are
  statements about two theories, and a model companion need not eliminate
  quantifiers; it need only be model complete.
\item A \emph{model completion} adds completeness over each atomic diagram
  of a model of the base theory. Elimination in the companion is a sufficient
  condition (\cref{thm:p03-w03-p03-020-qe-completion}), not the definition.
\end{enumerate}
\end{hbwarning}

\section{Back-and-forth proofs of quantifier elimination}
\label{sec:p03-w03-p03-020-back-forth}

\begin{hbdefinition}{def:p03-w03-p03-020-partial-iso}{Finite partial isomorphisms and back-and-forth systems}
Let \(\mathcal M\) and \(\mathcal N\) be \(\Lang\)-structures. A
\emph{finite partial isomorphism} from \(\mathcal M\) to \(\mathcal N\) is
a map \(p\) from a finite subset of \(M\) to \(N\) such that, for every
atomic formula \(\alpha(x_1,\ldots,x_k)\) and every tuple \(\bar a\) of
elements of the domain of \(p\),
\(\mathcal M\entails\alpha(\bar a)\) iff
\(\mathcal N\entails\alpha(p\bar a)\). The empty map is included; it is a
finite partial isomorphism exactly when \(\mathcal M\) and \(\mathcal N\)
satisfy the same atomic sentences. A \emph{back-and-forth system} is a
nonempty set \(I\) of finite partial isomorphisms such that for each
\(p\in I\): for each \(c\in M\) some \(q\in I\) extends \(p\) and has \(c\)
in its domain (\emph{forth}), and for each \(d\in N\) some \(q\in I\)
extends \(p\) and has \(d\) in its range (\emph{back}).
\end{hbdefinition}

Here atomic formulas may contain arbitrary terms, so a partial isomorphism
need not have a domain closed under the function symbols. In a relational
signature this is Fra\"iss\'e's notion of
\cref{def:p03-w03-p03-038-bf-system}, without its bound on the number of
rounds.

\begin{hbtheorem}{thm:p03-w03-p03-020-bf-preserve}{proved=annex:p03-w03-p03-020-bf-preserve}{Back-and-forth systems preserve all formulas}
Let \(I\) be a back-and-forth system between \(\mathcal M\) and
\(\mathcal N\). For every \(p\in I\), every formula \(\varphi(\bar x)\) and
every tuple \(\bar a\) of elements of the domain of \(p\),
\(\mathcal M\entails\varphi(\bar a)\) iff
\(\mathcal N\entails\varphi(p\bar a)\).
\end{hbtheorem}

\begin{hbcorollary}{thm:p03-w03-p03-020-bf-criterion}{proved=annex:p03-w03-p03-020-bf-criterion}{Back-and-forth criterion for elimination}
Suppose that for all models \(\mathcal M,\mathcal N\) of \(T\), every
finite partial isomorphism from \(\mathcal M\) to \(\mathcal N\) can be
extended to include any given element of \(M\) in its domain and, likewise,
any given element of \(N\) in its range. Then \(T\) eliminates
quantifiers, and any two models of \(T\) satisfying the same atomic
sentences are elementarily equivalent.
\end{hbcorollary}

\begin{hbtheorem}{thm:p03-w03-p03-020-dlo-bf}{proved=annex:p03-w03-p03-020-dlo-bf}{Elimination for dense linear orders by back-and-forth}
Let \(\mathcal M,\mathcal N\) be models of \(\mathsf{DLO}\) in \(\{<\}\).
Every finite partial isomorphism from \(\mathcal M\) to \(\mathcal N\)
extends to include any given element of \(M\) in its domain, and any given
element of \(N\) in its range. Consequently \(\mathsf{DLO}\) eliminates
quantifiers, is complete, and is model complete.
\end{hbtheorem}

This proves the cited \cref{thm:dlo-qe} and the corollary that follows it in
\cref{sec:qe-categoricity-01}. The back-and-forth proof yields no formula;
the explicit elimination step of \cref{def:p03-w03-p03-022-dlo-step},
proved correct in \cref{thm:p03-w03-p03-022-dlo-step}, produces the
quantifier-free equivalent and covers every combination of literals. The
inconsistent combinations reduce to \(\bot\) there, either directly
(\(y<y\)) or after the bound comparison produces an irreflexive atom, as in
the next example.

\begin{hbexample}{def:p03-w03-p03-020-dlo-literals}{Consistent and inconsistent literal combinations}
Modulo \(\mathsf{DLO}\):
\begin{enumerate}
\item \(\exists y\,(x<y\land y<z)\) is equivalent to \(x<z\).
\item \(\exists y\,(x<y\land y<x)\) is equivalent to \(x<x\), hence to
  \(\bot\): the bound comparison yields an irreflexive atom.
\item \(\exists y\,(y=x\land y<z\land z<y)\) is equivalent, by substituting
  \(x\) for \(y\), to \(x<z\land z<x\), which is false in every linear order.
\item \(\exists y\,\neg(y<x)\) is equivalent to
  \(\exists y\,(x<y)\lor\exists y\,(y=x)\), hence to \(\top\).
\item \(\exists y\,(y\ne x\land y\ne z)\) becomes the disjunction of
  \(\exists y\,(y<x\land y<z)\), \(\exists y\,(x<y\land z<y)\),
  \(\exists y\,(z<y\land y<x)\) and \(\exists y\,(x<y\land y<z)\). The
  first two have witnesses because there are no endpoints, so the formula
  is equivalent to \(\top\), although the last two disjuncts alone reduce
  to \(z<x\) and \(x<z\).
\end{enumerate}
In the back-and-forth proof these cases do not need to be listed: a formula
that is false at every point of one model is false at every point of every
model, because the partial isomorphisms carry truth both ways.
\end{hbexample}

\section{Dense linear orders and dense systems}
\label{sec:p03-w03-p03-020-dense}

Dense linear orders come in four kinds according to the endpoints present.
Only the kind without endpoints eliminates quantifiers in \(\{<\}\). The
other three eliminate once their endpoints are named.

\begin{hbdefinition}{def:p03-w03-p03-020-dense-variants}{Dense linear orders with prescribed endpoints}
Let \(E\subseteq\{\mathsf{b},\mathsf{t}\}\) and let \(\Lang_E\) be
\(\{<\}\) together with a constant symbol for each member of \(E\). The
\(\Lang_E\)-theory \(\mathsf{DLO}^{E}\) says that \(<\) is a strict linear
order with at least two elements, that between any two elements there is a
third, that \(\mathsf{b}\) is the least element if
\(\mathsf{b}\in E\) and that there is no least element otherwise, and that
\(\mathsf{t}\) is the greatest element if \(\mathsf{t}\in E\) and that
there is no greatest element otherwise. Thus \(\mathsf{DLO}^{\emptyset}\)
is \(\mathsf{DLO}\) together with the axiom \(\exists x\,\exists y\,x\ne y\),
which \(\mathsf{DLO}\) already entails.
\end{hbdefinition}

\begin{hbtheorem}{thm:p03-w03-p03-020-dense-variants}{proved=annex:p03-w03-p03-020-dense-variants}{Elimination for each endpoint pattern}
For each \(E\subseteq\{\mathsf{b},\mathsf{t}\}\), the theory
\(\mathsf{DLO}^{E}\) eliminates quantifiers in \(\Lang_E\), is complete and
is model complete. Every linear order embeds in a model of
\(\mathsf{DLO}^{\emptyset}\), so \(\mathsf{DLO}\) is the model completion of
the theory of linear orders in \(\{<\}\).
\end{hbtheorem}

\begin{hbwarning}{def:p03-w03-p03-020-endpoint-failure}{Unnamed endpoints block elimination}
Let \(\mathsf{b}\in E\), and let \(T\) be the \(\{<\}\)-theory whose models
are the \(\{<\}\)-reducts of models of \(\mathsf{DLO}^{E}\). Then \(T\) is
complete but neither eliminates quantifiers nor is model complete. A
quantifier-free \(\{<\}\)-formula in one variable \(x\) is a Boolean
combination of \(x<x\) and \(x=x\), so it is true of all elements of a
model or of none; but \(\exists y\,(y<x)\) fails at the least element and
holds at every other element. And
\(\mathcal M=(\mathbb Q\cap[0,1];<)\subseteq\mathcal N=(\mathbb Q\cap[-1,1];<)\)
are models when \(E=\{\mathsf{b},\mathsf{t}\}\), but \(\exists y\,(y<0)\) is
true in \(\mathcal N\) and false in \(\mathcal M\). The same failure holds
with \(\mathsf{t}\) in place of \(\mathsf{b}\). The proof that \(T\) is
complete is part of the annex proof of
\cref{thm:p03-w03-p03-020-dense-variants}.
\end{hbwarning}

\begin{hbtheorem}{thm:p03-w03-p03-020-marked-point}{proved=annex:p03-w03-p03-020-marked-point}{A model complete theory without elimination}
Let \(Z\) be a unary relation symbol and let \(T_Z\) be the
\(\{<,Z\}\)-theory consisting of \(\mathsf{DLO}\) and the sentence saying
that exactly one element satisfies \(Z\). Then \(T_Z\) is model complete and
does not eliminate quantifiers.
\end{hbtheorem}

\section{Algebraically closed fields and real closed fields}
\label{sec:p03-w03-p03-020-fields}

The ring language is \(\Lang_{\mathrm r}=\{0,1,+,-,\cdot\}\), with \(-\)
unary. Every \(\Lang_{\mathrm r}\)-term in variables \(\bar x\) denotes, in
every commutative ring, the value of a polynomial with integer coefficients,
so every atomic formula is equivalent modulo the field axioms to one of the
form \(f(\bar x)=0\) with \(f\in\mathbb Z[\bar x]\)
\cite[Remark~2.26, p.~19]{Kruckman2020}. For \(p\) prime, \(\chi_p\) is the
atomic sentence \(1+\cdots+1=0\) with \(p\) summands. The theory
\(\mathsf{Fld}_p\) consists of the field axioms with \(\chi_p\) if \(p\) is
prime, and with \(\neg\chi_q\) for every prime \(q\) if \(p=0\). The theory
\(\mathsf{ACF}_p\) adds, for each \(d\ge1\), the sentence
\(\forall a_0\cdots\forall a_{d-1}\,\exists y\,
(y^d+a_{d-1}y^{d-1}+\cdots+a_0=0)\). The characteristic is fixed by the
theory; without it the class of algebraically closed fields is not complete,
as \cref{tab:p03-qe-catalog} records.

\begin{hblemma}{thm:p03-w03-p03-020-field-algebra}{proved=annex:p03-w03-p03-020-field-algebra}{Field extensions used in the elimination}
Let \(E\subseteq K\) be fields and \(\sigma:E\to L\) a field embedding.
\begin{enumerate}
\item If \(g\in E[X]\) is irreducible, then \(E[X]/(g)\) is a field
  extension of (a copy of) \(E\) in which the class of \(X\) is a root of
  \(g\). Hence finitely many nonconstant polynomials over \(E\) all have
  roots in one extension field of \(E\).
\item If \(\alpha\in K\) is algebraic over \(E\) with minimal polynomial
  \(g\), and \(\beta\in L\) is a root of \(\sigma g\), then \(\sigma\)
  extends to an embedding \(E(\alpha)\to L\) with \(\alpha\mapsto\beta\).
  If \(L\) is algebraically closed such a \(\beta\) exists.
\item If \(\alpha\in K\) is transcendental over \(E\) and \(\beta\in L\)
  is transcendental over \(\sigma(E)\), then \(\sigma\) extends to an
  embedding \(E(\alpha)\to L\) with \(\alpha\mapsto\beta\).
\item An injective ring homomorphism from an integral domain \(A\) into a
  field \(L\) extends uniquely to an embedding of the fraction field of
  \(A\) into \(L\).
\item Every algebraically closed field is infinite.
\end{enumerate}
\end{hblemma}

\begin{hbtheorem}{thm:p03-w03-p03-020-closure-exists}{proved=annex:p03-w03-p03-020-closure-exists}{Algebraically closed extensions exist}
Assume the Axiom of Choice. Every field \(K\) is a subfield of an
algebraically closed field of the same characteristic, and the elements of
that field algebraic over \(K\) form an algebraically closed field
algebraic over \(K\).
\end{hbtheorem}

\begin{hbtheorem}{thm:p03-w03-p03-020-acf-qe}{proved=annex:p03-w03-p03-020-acf-qe}{Elimination for algebraically closed fields of fixed characteristic}
For \(p=0\) and for every prime \(p\), \(\mathsf{ACF}_p\) eliminates
quantifiers in \(\Lang_{\mathrm r}\). It is complete and model complete.
\end{hbtheorem}

This proves the cited theorem on algebraically closed fields in
\cref{sec:qe-categoricity-01}. Kruckman's proof uses the same test and takes
the existence of algebraic closures and the extension of embeddings as known
\cite[Theorem~2.27, p.~20]{Kruckman2020}; here they are
\cref{thm:p03-w03-p03-020-field-algebra,thm:p03-w03-p03-020-closure-exists}.

\begin{hbtheorem}{thm:p03-w03-p03-020-acf-completion}{proved=annex:p03-w03-p03-020-acf-completion}{\(\mathsf{ACF}_p\) is the model completion of \(\mathsf{Fld}_p\)}
Assume the Axiom of Choice. For \(p=0\) and every prime \(p\):
\begin{enumerate}
\item every model of \(\mathsf{Fld}_p\) embeds in a model of
  \(\mathsf{ACF}_p\), so \(\mathsf{ACF}_p\) and \(\mathsf{Fld}_p\) have
  the same universal consequences;
\item \(\mathsf{ACF}_p\) is model complete;
\item for every field \(K\) of characteristic \(p\), the
  \(\Lang_{\mathrm r}(K)\)-theory
  \(\mathsf{ACF}_p\cup\operatorname{Diag}(K)\) is consistent and complete;
\item the fields of characteristic \(p\) that are existentially closed
  among fields of characteristic \(p\) are exactly the algebraically closed
  ones.
\end{enumerate}
\end{hbtheorem}

This proves the \(\mathsf{ACF}_p\) half of the cited
\cref{thm:p03-field-model-completions}.

\begin{hbexample}{def:p03-w03-p03-020-quadratic}{Eliminating a root quantifier}
Modulo \(\mathsf{ACF}_p\), the formula \(\exists y\,(ay^2+by+c=0)\) is
equivalent to \(a\ne0\lor b\ne0\lor c=0\). If \(a\ne0\), the monic
polynomial \(y^2+(b/a)y+c/a\) has a root by the axiom for \(d=2\); if
\(a=0\) and \(b\ne0\), then \(y=-c/b\) is a root; if \(a=b=0\), the
equation reads \(c=0\). The equivalence fails in the theory of fields:
\(\mathbb R\subseteq\mathbb C\), and \(y^2+1=0\) has a root in
\(\mathbb C\) but not in \(\mathbb R\), so
\(\exists y\,(y^2+1=0)\) is not preserved down this embedding and is not
equivalent to a quantifier-free formula modulo the field axioms
\cite[Example~2.21, p.~16]{Kruckman2020}.
\end{hbexample}

\begin{hbwarning}{def:p03-w03-p03-020-language-choice}{Characteristic, order and definable choice}
Three hypotheses stay attached to these results. The characteristic must
be fixed for completeness; the class of all algebraically closed fields
passes the one-point test by the same proof, since equal quantifier-free
types already agree on each \(\chi_p\), but it leaves each \(\chi_p\)
undecided. For the real field, elimination is possible only in a language
that contains the order: in \(\Lang_{\mathrm r}\) the formula
\(\exists y\,(y^2=x)\) defines the nonnegative elements of \(\mathbb R\),
which is neither finite nor cofinite, whereas every quantifier-free
\(\Lang_{\mathrm r}\)-formula in one variable defines, in any field, a
finite or cofinite set, because an atom \(f(x)=0\) with
\(f\in\mathbb Z[x]\) defines the whole field when \(f\) is zero in the
field and at most \(\deg f\) points otherwise, and finite and cofinite
sets are closed under Boolean operations. Finally, elimination does not
by itself provide definable choice of witnesses
(\cref{def:p03-definable-skolem-functions}): the conjugate roots in
\(\mathsf{ACF}_p\) cannot be selected uniformly, while the catalog
\cref{tab:p03-qe-catalog} records definable choice for real closed fields
as a separate property.
\end{hbwarning}


\subsection{Ordered fields, real closures and real closed fields}
\label{sec:p03-w03-p03-020-rcf}

The ordered analogue of the passage from fields to algebraically closed
fields runs in the ordered-ring language
\(\Lang_{\mathrm{or}}=\{0,1,+,-,\mathbin{\cdot},<\}\) of
\cref{def:p03-w03-p03-037-rcf}, in which \(-\) is binary subtraction rather
than the unary minus of \(\Lang_{\mathrm r}\). Terms again denote integer
polynomials (\cref{def:p03-w03-p03-037-parametric}), but the atomic formulas
now include strict inequalities between polynomials as well as equations.
The inputs are an ordered field and the axioms of \(\mathsf{RCF}\); the
conclusions are that every ordered field has a real closure, that
\(\mathsf{RCF}\) is the model completion of the theory of ordered fields,
with completeness over the \(\Lang_{\mathrm{or}}\)-diagram of each ordered
field, and that the existentially closed ordered fields are exactly the real
closed fields. Elimination of quantifiers for \(\mathsf{RCF}\) in
\(\Lang_{\mathrm{or}}\) is \cref{thm:p03-w03-p03-037-qe}, proved from the
axioms in \cref{sec:p03-w03-p03-037-qe}; it is used here and not reproved.

\begin{hbdefinition}{def:p03-w03-p03-020-ordered-fields}{Ordered fields and real closures}
\(\mathsf{OF}\) is the \(\Lang_{\mathrm{or}}\)-theory of ordered fields of
\cref{def:p03-w03-p03-037-rcf}: the axioms of commutative fields (with
\(\forall x\,\exists y\,(x=0\lor x\mathbin{\cdot}y=1)\) for inverses), the
axioms of a strict linear order for \(<\), and the two compatibility axioms
\(x<y\rightarrow x+z<y+z\) and \((0<x\land0<y)\rightarrow0<x\mathbin{\cdot}y\).
\(\mathsf{RCF}\) is \(\mathsf{OF}\) together with the intermediate-value
axioms of \cref{def:p03-w03-p03-037-rcf}. For ordered fields \((F,<)\) and
\((L,<_L)\) with \(F\subseteq L\), the order \(<_L\) \emph{extends} \(<\)
when \((F,<)\) is an \(\Lang_{\mathrm{or}}\)-substructure of \((L,<_L)\),
that is, \(F\) is a subfield of \(L\) and \(<\) is the restriction of
\(<_L\) to \(F\). A \emph{real closure} of an ordered field \((F,<)\) is an
ordered field \((R,<_R)\models\mathsf{RCF}\) such that \(R\) is a field
extension of \(F\), algebraic over \(F\), and \(<_R\) extends \(<\).
\index{real closure}
\index{ordered field}
\end{hbdefinition}

The order is part of the data. Kruckman defines a real closure of a field
without an order, as an algebraic extension that is formally real and has no
proper formally real algebraic extension \cite[Definitions~3.16 and~3.18,
p.~30]{Kruckman2020}, and then shows that for an ordered field it can be
chosen so that its unique order extends the given one
\cite[Theorem~3.28, p.~36]{Kruckman2020}. The definition above builds that
requirement in. For the agreement of his definition of real closed fields
with the intermediate-value form used here see
\cite[Theorem~3.31, p.~37]{Kruckman2020};
\cref{thm:p03-w03-p03-020-rcf-char} proves the part used below.

\begin{hblemma}{thm:p03-w03-p03-020-order-extension}{proved=annex:p03-w03-p03-020-order-extension}{Extending an order}
Assume the Axiom of Choice. Let \((K,<)\) be an ordered field and
\(L\supseteq K\) a field. The order \(<\) extends to an order \(<_L\) making
\((L,<_L)\) an ordered field if and only if there are no \(m\ge1\), no
\(p_1,\ldots,p_m\in K\) with \(0<p_i\), and no \(y_1,\ldots,y_m\in L\) such
that
\[
 -1=p_1y_1^2+\cdots+p_my_m^2 .
\]
\end{hblemma}

This is the positive-cone argument of \cite[Lemmas~3.12 and~3.13,
pp.~27--28]{Kruckman2020}, run over the positive elements of \(K\) instead
of over sums of squares alone, so that the order obtained on \(L\) contains
the given order on \(K\).

\begin{hblemma}{thm:p03-w03-p03-020-root-ext}{proved=annex:p03-w03-p03-020-root-ext}{Adjoined roots keep an order}
Assume the Axiom of Choice. Let \((K,<)\) be an ordered field, \(L\supseteq
K\) a field and \(\alpha\in L\).
\begin{enumerate}
\item If \(\alpha^2=a\) for some \(a\in K\) with \(0<a\), then \(<\)
  extends to an order of the field \(K(\alpha)\).
\item If \(\alpha\) is a root of an irreducible polynomial \(g\in K[X]\) of
  odd degree, then \(<\) extends to an order of \(K(\alpha)\).
\end{enumerate}
\end{hblemma}

Kruckman proves the unordered forms, for formally real fields
\cite[Lemmas~3.20 and~3.21, p.~31]{Kruckman2020}; the descent on the degree
in item~2 is the same argument.

\begin{hbtheorem}{thm:p03-w03-p03-020-rcf-char}{proved=annex:p03-w03-p03-020-rcf-char}{Real closed fields by roots}
For an ordered field \((R,<)\) the following are equivalent.
\begin{enumerate}
\item \((R,<)\models\mathsf{RCF}\).
\item Every \(a\in R\) with \(0<a\) is a square in \(R\), and every
  polynomial of odd degree in \(R[X]\) has a root in \(R\).
\end{enumerate}
When they hold, \(-1\) is not a square in \(R\), and the field
\(R(\mathrm i)=R[X]/(X^2+1)\), with \(\mathrm i\) the class of \(X\), is
algebraically closed.
\end{hbtheorem}

The implication from 2 to 1 is where algebra enters. The annex proof first
shows that \(R(\mathrm i)\) is algebraically closed, by induction on the
power of \(2\) dividing the degree and with symmetric polynomials. It then
factors a polynomial over \(R\) into linear factors and quadratic factors
without roots, and a sign change between two points must come from a linear
factor. Kruckman proves that \(R(\mathrm i)\) is algebraically closed with
Galois theory and Sylow subgroups \cite[Theorem~3.22,
pp.~32--33]{Kruckman2020} and then derives the intermediate value property
\cite[Theorem~3.31, pp.~37--38]{Kruckman2020}.

\begin{hbtheorem}{thm:p03-w03-p03-020-real-closure}{proved=annex:p03-w03-p03-020-real-closure}{Real closures exist}
Assume the Axiom of Choice. Let \((F,<)\) be an ordered field and let
\(\Omega\supseteq F\) be an algebraically closed field algebraic over \(F\),
as provided by \cref{thm:p03-w03-p03-020-closure-exists}. Then some subfield
\(R\) of \(\Omega\) containing \(F\) carries an order \(<_R\) such that
\((R,<_R)\) is a real closure of \((F,<)\). In particular every ordered field
has a real closure.
\end{hbtheorem}

Kruckman first adjoins square roots of all positive elements and then takes
a maximal formally real algebraic extension
\cite[Lemma~3.27 and Theorem~3.28, p.~36]{Kruckman2020}; the annex proof
instead takes a maximal ordered subfield of \(\Omega\) whose order extends
\(<\), so that the extension of the order is kept at every stage. Real
closures are also unique up to isomorphism over \(F\), by Sturm's theorem
\cite[Theorem~3.38, p.~42]{Kruckman2020}; that uniqueness is not proved here
and nothing below uses it. Item~3 of the next theorem gives the weaker
statement that any two real closures of \((F,<)\) satisfy the same
\(\Lang_{\mathrm{or}}(F)\)-sentences.

\begin{hbtheorem}{thm:p03-w03-p03-020-rcf-mc}{proved=annex:p03-w03-p03-020-rcf-mc}{Model completion of ordered fields}
Assume the Axiom of Choice. In the ordered-ring language
\(\Lang_{\mathrm{or}}\):
\begin{enumerate}
\item every ordered field embeds in a model of \(\mathsf{RCF}\), so
  \(\mathsf{OF}\) and \(\mathsf{RCF}\) have the same universal consequences;
\item \(\mathsf{RCF}\) eliminates quantifiers, and it is complete and model
  complete;
\item for every ordered field \((F,<)\), the
  \(\Lang_{\mathrm{or}}(F)\)-theory
  \(\mathsf{RCF}\cup\operatorname{Diag}(F,<)\) is consistent and complete,
  so \(\mathsf{RCF}\) is the model completion of \(\mathsf{OF}\); in
  particular any two real closures of \((F,<)\), with each constant
  \(c_a\) naming \(a\in F\), satisfy the same
  \(\Lang_{\mathrm{or}}(F)\)-sentences;
\item the ordered fields that are existentially closed among ordered
  fields, in the sense of \cref{def:p03-existentially-closed}, are exactly
  the real closed fields.
\end{enumerate}
\end{hbtheorem}

This proves the real-closed-field half of the cited
\cref{thm:p03-field-model-completions}, in the ordered-ring language; the
algebraically closed half is \cref{thm:p03-w03-p03-020-acf-completion}.
Kruckman proves elimination, completeness and model completeness
\cite[Theorem~3.39, p.~43; Corollary~3.40, p.~44; Theorem~3.46,
p.~46]{Kruckman2020}, and his elimination proof uses the uniqueness of real
closures; the route here goes through \cref{thm:p03-w03-p03-037-qe}, which
does not need it.

\begin{hbexample}{def:p03-w03-p03-020-rcf-example}{Deciding sentences over an ordered field}
\emph{The ordered field.} In \(\mathbb Q(t)\), with \(t\) transcendental
over \(\mathbb Q\), call a nonzero element \(f/g\) (\(f,g\in\mathbb Q[t]\)
nonzero) positive when the leading coefficients satisfy
\(\operatorname{lc}(f)\operatorname{lc}(g)>0\). This does not depend on the
representation: \(f/g=f'/g'\) gives
\(\operatorname{lc}(f)\operatorname{lc}(g')=\operatorname{lc}(f')\operatorname{lc}(g)\),
and multiplying by \(\operatorname{lc}(g)\operatorname{lc}(g')\) shows that
\(\operatorname{lc}(f)\operatorname{lc}(g)\) and
\(\operatorname{lc}(f')\operatorname{lc}(g')\) have the same sign. Write
positive elements with denominators of positive leading coefficient. Then
products of positive elements are positive, and so are sums, because the
leading coefficient of \(fg'+f'g\) is that of the summand of larger degree,
or the sum of two positive rationals when the degrees agree. Exactly one of
\(x\) and \(-x\) is positive for \(x\ne0\). So \(x<_+y\) iff \(y-x\) is
positive defines an ordered field \((\mathbb Q(t),<_+)\) extending the order
of \(\mathbb Q\), in which \(n<_+t\) for every integer \(n\), because
\(t-n\) has leading coefficient \(1\).

\emph{One quantifier, eliminated.} Modulo \(\mathsf{RCF}\),
\(\exists y\,(y\mathbin{\cdot}y=x)\) is equivalent to \(0<x\lor x=0\):
squares are nonnegative in any ordered field, and nonnegative elements of a
real closed field are squares by
\cref{thm:p03-w03-p03-020-rcf-char}. So for every
\(a\in\mathbb Q(t)\), the theory
\(\mathsf{RCF}\cup\operatorname{Diag}(\mathbb Q(t),<_+)\) entails
\(\exists y\,(y\mathbin{\cdot}y=c_a)\) when \(a\) is positive or zero,
because the atomic sentence \(c_0<c_a\) or \(c_a=c_0\) is then in the
diagram, and entails its negation otherwise. For instance it entails
\(\exists y\,(y\mathbin{\cdot}y=c_{t-n})\) for every integer \(n\), so every
real closure of \((\mathbb Q(t),<_+)\) contains square roots of all the
elements \(t-n\). None of them lies in \(\mathbb Q(t)\): an equation
\(f^2=(t-n)g^2\) with nonzero \(f,g\in\mathbb Q[t]\) is impossible, since
the left side has even degree and the right side odd degree.

\emph{Existential closure fails below.} The ordered field
\((\mathbb Q,<)\) is not existentially closed among ordered fields: the
existential sentence \(\exists y\,(y\mathbin{\cdot}y=1+1)\) is false in
\(\mathbb Q\) and true in every real closure of \(\mathbb Q\), which exists by
\cref{thm:p03-w03-p03-020-real-closure}. This is the kind of witness used in
the proof of item~4 of \cref{thm:p03-w03-p03-020-rcf-mc}.
\end{hbexample}

\begin{hbwarning}{def:p03-w03-p03-020-rcf-ring-diagram}{The order belongs in the diagram}
Completeness over diagrams in \cref{thm:p03-w03-p03-020-rcf-mc}
needs the order of \(F\) in the diagram. Let
\(\sigma:\mathbb Q(t)\to\mathbb Q(t)\) be the field automorphism with
\(\sigma(t)=-t\) (the substitution \(t\mapsto-t\) on \(\mathbb Q[t]\), an
injective ring homomorphism into \(\mathbb Q(t)\), extended to fractions by
item~4 of \cref{thm:p03-w03-p03-020-field-algebra}; it is its own inverse),
and put \(x<_-y\) iff \(\sigma(x)<_+\sigma(y)\), with \(<_+\) as in
\cref{def:p03-w03-p03-020-rcf-example}. Since \(\sigma\) preserves sums and
products, \((\mathbb Q(t),<_-)\) is an ordered field, and \(t<_-0\) because
\(\sigma(t)=-t<_+0\). Let \(R_+\) and \(R_-\) be real closures of
\((\mathbb Q(t),<_+)\) and \((\mathbb Q(t),<_-)\)
(\cref{thm:p03-w03-p03-020-real-closure}). As
\(\Lang_{\mathrm r}(\mathbb Q(t))\)-structures, both satisfy the
\(\Lang_{\mathrm r}\)-sentences true in all real closed fields and the
\(\Lang_{\mathrm r}\)-diagram of the field \(\mathbb Q(t)\). But
\(R_+\models\exists y\,(y\mathbin{\cdot}y=c_t)\), since \(0<t\) there, and
\(R_-\models\neg\exists y\,(y\mathbin{\cdot}y=c_t)\), since \(t<0\) there and
squares are nonnegative. So the ring diagram of a field with more than one
order does not decide the theory of its real closures, and real closures
depend on the chosen order, as Kruckman observes for
\(\mathbb Q(\sqrt2)\) \cite[Example~3.4, p.~25]{Kruckman2020}. Elimination
also needs the order symbol: in \(\Lang_{\mathrm r}\) it fails for the real
closed fields (\cref{thm:p03-w03-p03-037-ring-language} and
\cref{def:p03-w03-p03-020-language-choice}).
\end{hbwarning}
